\begin{corollary}[3.3.3]
\label{IV.2.3.3.3-fr}
Bajo las hipótesis de \textup{(3.3.1)}, supongamos además que $X$ e $Y$ son localmente
noetherianos, y que $\mathcal E$ y $\mathcal F$ son coherentes. Entonces las
condiciones siguientes son equivalentes:
\begin{enumerate}
\item[\emph{a)}] $\mathcal E\otimes_Y\mathcal F$ carece de ciclo primo
asociado inmerso.
\item[\emph{b)}] Para todo punto
$y\in\operatorname{Ass}(\mathcal E)\cap f(\operatorname{Supp}(\mathcal F))$,
$\overline{\{y\}}$ es un ciclo primo asociado no inmerso de $\mathcal E$ y
$\mathcal F_y$ carece de ciclo primo asociado inmerso.
\end{enumerate}
\end{corollary}

\oldpage[IV-2]{44}
Supongamos que se verifica \emph{a)}. Las hipótesis implican que
$\mathcal E\otimes_Y\mathcal F$ es un $\mathcal O_X$-módulo coherente
$(\mathbf{0}_{\mathrm I},5.3.11\text{ y }5.3.5)$; sus ciclos primos asociados son,
por tanto, las componentes irreducibles de
$\operatorname{Supp}(\mathcal E\otimes_Y\mathcal F)$ (3.1.4) y, para todo
punto maximal $x$ de $\operatorname{Supp}(\mathcal E\otimes_Y\mathcal F)$,
$f(x)=y$ es un punto maximal de $\operatorname{Supp}(\mathcal E)$ y $x$ un
punto maximal de $\operatorname{Supp}(\mathcal F_y)$ (2.5.5). Como, en virtud
de (3.3.1) y del hecho de que la relación
$y\in f(\operatorname{Supp}(\mathcal F))$ implica $\mathcal F_y\ne0$
(\textbf I, 9.1.13), todo punto de
$\operatorname{Ass}(\mathcal E)\cap f(\operatorname{Supp}(\mathcal F))$ es la
imagen por $f$ de un punto maximal de
$\operatorname{Supp}(\mathcal E\otimes_Y\mathcal F)$, se ve que se verifica la
condición \emph{b)}.

Recíprocamente, supongamos que se verifica \emph{b)} y demostremos que, si
$z$, $z'$ son dos puntos distintos de
$\operatorname{Ass}(\mathcal E\otimes_Y\mathcal F)$, ninguno de ellos puede
ser adherente al otro. En primer lugar, si $f(z)=f(z')=y$, se tiene
$y\in\operatorname{Ass}(\mathcal E)$ y $z$ y $z'$ pertenecen a
$\operatorname{Ass}(\mathcal F_y)$ por (3.3.1), de donde
$y\in f(\operatorname{Supp}(\mathcal F))$; como, por hipótesis, ninguno de los
dos puntos $z$, $z'$ es adherente al otro \emph{en} $f^{-1}(y)$, ninguno de
ellos puede ser adherente al otro \emph{en} $X$. Si $y=f(z)$ e $y'=f(z')$ son
distintos, pertenecen a
$\operatorname{Ass}(\mathcal E)\cap f(\operatorname{Supp}(\mathcal F))$, de
modo que ninguno puede ser adherente al otro en $Y$; de la continuidad de $f$
resulta que ninguno de los puntos $z$, $z'$ puede ser adherente al otro en $X$.

